% ===================================================================
%  CUADRO N.º 1 — La Escuela. — El Instituto. — El Aula.
%
%  Ceci n'est PAS une transcription : c'est une traduction, faite en
%  2026, en espagnol d'aujourd'hui. D'ou trois differences avec
%  texto/io et texto/fr, et trois seulement :
%
%   * pas de \nl ni de \cc — la traduction ne suit aucune ligne
%     imprimee, et les lignes de ce fichier ne sont que du confort ;
%   * pas de \begin{VUpage} — elle n'a ni feuillet ni folio, et la
%     page de lecture ne lui en affiche pas ;
%   * le gras se pose sur le TERME seul, ponctuation dehors.
%
%  Tout le reste est identique, et doit l'etre : les cles « %%K » sont
%  celles de texto/io, macro pour macro pour l'apparat, et les renvois
%  \textsuperscript{(n)} visent les MEMES objets de la planche — ce
%  sont eux qui ouvrent les gros plans.
%
%  L'ido est le texte de base : c'est lui que la colonne de gauche
%  porte, et c'est donc lui qu'on traduit. Le francais et l'anglais
%  servent de controle ; l'anglais surtout, qui a deja tranche chaque
%  question de vocabulaire, et sur lequel l'espagnol se cale mot pour
%  mot afin que les deux traductions disent la meme chose.
%
%  LA FAUTE N'ETAIT PAS DANS L'ESPAGNOL, ELLE ETAIT DANS LA MANIERE DE
%  TRADUIRE. Vingt-deux alineas de cette colonne ouvraient par « he
%  aqui », decalque du « voici » que le fac-simile francais place en
%  tete de presque chaque objet nouveau. La tournure existe en
%  espagnol, mais elle y est litteraire : personne ne l'ecrit en 2026
%  pour dire ou se tient la laitiere. On l'a remplacee cas par cas —
%  « a su lado esta la lechera », « en la mesa hay » — jamais par une
%  formule unique, puisque le francais ne disait pas non plus la meme
%  chose vingt-deux fois.
%
%  La preuve que la faute vient de la source et non de la langue : les
%  MEMES blocs portaient « ty » en suedois, « for » en anglais, « eis »
%  en portugais. Une colonne se relit donc contre le fac-simile, pas
%  contre elle-meme.
% ===================================================================

%%K t01-apar-0 apar
\VUsaut{2.0mm}
\VUcentre{12.6pt}{120}{FOLLETO EXPLICATIVO}
\VUsaut{3.2mm}
\VUcentre{8.4pt}{160}{DE}
\VUsaut{2.6mm}
\VUtitre{15.6pt}{0}{los Cuadros Auxiliares Delmas}
\VUsaut{3.4mm}
\VUornamento{9pt}
\VUsaut{5.0mm}
\VUcentre{13.2pt}{90}{PRIMERA SERIE}
\VUsaut{3.0mm}
\VUfilet{16mm}
\VUsaut{5.6mm}

%%K t01-apar-1 apar
\VUtitre{15.0pt}{60}{CUADRO N.º 1}
\VUsaut{1.4mm}
\VUtitre{11.4pt}{0}{La Escuela. --- El Instituto. --- El Aula.}
\VUsaut{2.2mm}
\VUfilet{20mm}
\VUsaut{6.4mm}

%%K t01-tit-1 sub
\VUpk{10.2pt}{40}{Descripción general.}
\VUsaut{3.0mm}

%%K t01-01-1 p
1. --- Este cuadro representa un \VUgras{aula} de \VUgras{instituto}. En esta
aula hay ocho \VUgras{mesas} \textsuperscript{(18)} y ocho \VUgras{bancos}
\textsuperscript{(32)}. En cada banco se sientan tres alumnos. El
\VUgras{profesor} \textsuperscript{(7)} está sentado en una \VUgras{silla}
\textsuperscript{(8)} en su \VUgras{cátedra} \textsuperscript{(6)}.

%%K t01-02-1 p
2. --- A la izquierda de la cátedra hay un \VUgras{armario} con puertas de
cristal \textsuperscript{(15)}. A la derecha está la \VUgras{pizarra}
\textsuperscript{(1)} con el \VUgras{borrador} \textsuperscript{(2)} y la \VUgras{tiza}
\textsuperscript{(3)}.

%%K t01-02-2 p
Encima de la cátedra cuelgan unas \VUgras{láminas murales} \textsuperscript{(9, 11,
12)} y un \VUgras{mapa} \textsuperscript{(10)}.

%%K t01-03-1 p
3. --- No todos los alumnos están en su \VUgras{sitio} \textsuperscript{(17)}. Juan
\textsuperscript{(4)} está de pie junto a la pizarra; sostiene el borrador y borra
las \VUgras{figuras de geometría}.

%%K t01-03-2 p
Pedro está cerca de la cátedra; recoge los \VUgras{ejercicios escritos}
\textsuperscript{(5)} y se los lleva al profesor.

%%K t01-04-1 p
4. --- \VUgras{Este} profesor es el profesor de \VUgras{lenguas modernas}.
Enseña a los alumnos a hablar, leer y escribir lenguas extranjeras
\VUgras{(alemán, inglés, español, italiano, ruso} o \VUgras{francés)}.

%%K t01-05-1 p
5. --- El aula es muy amplia y muy clara. Al fondo hay una \VUgras{ventana}
ancha con dos \VUgras{montantes} \textsuperscript{(78)} y una \VUgras{puerta
acristalada} para airearla e iluminarla.

%%K t01-05-2 p
Esta aula no es fría. En medio, para calentarla, hay una \VUgras{estufa}
\textsuperscript{(46)} muy buena con su \VUgras{tubo} \textsuperscript{(45)}.

%%K t01-05-3 p
En el tabique hay fijadas unas \VUgras{perchas} \textsuperscript{(58)}; de ellas
cuelgan los gorros y los abrigos de los alumnos.

%%K t01-tit-2 sub
\VUsaut{5.2mm}
\VUpk{10.2pt}{40}{El patio de recreo.}
\VUsaut{3.6mm}

%%K t01-06-1 p
6. --- El \VUgras{reloj} \textsuperscript{(63)} marca las nueve y cinco.

%%K t01-06-2 p
El \VUgras{recreo} \textsuperscript{(76)} acaba de terminar y la clase empieza otra
vez. El recreo se hace en el \VUgras{patio} \textsuperscript{(75)}. Vemos a unos
cuantos alumnos que juegan. Un \VUgras{pasante} \textsuperscript{(74)} los vigila.

%%K t01-07-1 p
7. --- En el patio vemos además a otras personas: el \VUgras{inspector}
\textsuperscript{(73)}, el \VUgras{director} \textsuperscript{(71)} y el
\VUgras{jefe de estudios} \textsuperscript{(72)}.

%%K t01-07-2 p
El inspector inspecciona las escuelas, el director dirige el instituto,
el jefe de estudios vigila los estudios.

%%K t01-07-3 p
El cuarto personaje es el \VUgras{conserje} \textsuperscript{(77)}. Lleva bajo el
brazo un registro para anotar a los ausentes.

%%K t01-08-1 p
8. --- Al fondo del patio están los \VUgras{aparatos de gimnasia}
\textsuperscript{(70)} y el taller de \VUgras{trabajos manuales} \textsuperscript{(69)}.

%%K t01-08-2 p
A la derecha del cuadro alcanzamos a ver las \VUgras{aulas} de
\VUgras{música} y de \VUgras{dibujo}.

%%K t01-noto-1 noto
\VUnotes{143.2mm}{%
(*) Para los \textit{nombres de pila} se aconseja transcribirlos también
en su forma latina. Es el único modo de evitar variantes sin cuento. Por
lo demás, ¿cómo elegir entre \textit{Giovanni, Jean, Johann, Jan, Hans,}
\textit{John, Ivan,} etc.? ¿Habrá que crear un diccionario especial de
nombres de pila? Tomemos del latín su nominativo y digamos:
\VUgras{Ioannes, Ioanna; Iulius, Iulia; Iakobus; Andreas; Lukas.}
(Gramática completa).%
}

%%K t01-tit-3 sub
\VUpk{10.2pt}{40}{Descripción detallada.}
\VUsaut{2.4mm}

%%K t01-tit-4 sub
\VUpk{10.2pt}{40}{Los buenos alumnos.}
\VUsaut{3.0mm}

%%K t01-09-1 p
9. --- Los alumnos del primer banco a la derecha de la cátedra son
\VUgras{Enrique} (*), \VUgras{Esteban} y \VUgras{Carlos}.

%%K t01-09-2 p
Enrique es muy atento y trabajador; escucha al profesor. Sobre su mesa
tiene un \VUgras{cuaderno} \textsuperscript{(28)}, un \VUgras{libro} \textsuperscript{(29)},
un \VUgras{diccionario} \textsuperscript{(30)} y un \VUgras{plumier} \textsuperscript{(31)}.
Su libro tiene muchas \VUgras{páginas}; cada capítulo contiene varios
\VUgras{párrafos} y cada \VUgras{línea} muchas \VUgras{palabras}.

%%K t01-09-4 p
Su vecino Esteban \textsuperscript{(24)} es aplicado. Apunta en su cuaderno la
lección para la próxima vez. Tiene buena \VUgras{letra}. Su libro está
cerrado. Solo vemos la \VUgras{cubierta} \textsuperscript{(27)}. A su lado está su
\VUgras{compás} \textsuperscript{(26)}.

%%K t01-09-5 p
Carlos \textsuperscript{(23)} sostiene su libro en la mano; lo abre para repasar
la lección. Como todos los alumnos, tiene delante un \VUgras{tintero}
\textsuperscript{(25)}. Es obediente.

%%K t01-10-1 p
10. --- En la mesa de al lado están \VUgras{Luis, Antonio} y \VUgras{Pablo}.
Luis recita su \VUgras{lección} \textsuperscript{(22)} y responde a las
\VUgras{preguntas} del profesor. Antonio \textsuperscript{(21)} está muy distraído;
a veces el profesor hasta lo castiga. Pablo \textsuperscript{(20)} es buen
compañero, amable y servicial. Es muy atento.

%%K t01-tit-5 sub
\VUsaut{4.2mm}
\VUpk{10.2pt}{40}{Los malos alumnos.}
\VUsaut{3.2mm}

%%K t01-11-1 p
11. --- Detrás de Enrique está \VUgras{Jorge} \textsuperscript{(38)}. No es mal
chico, pero es \VUgras{hablador}, distraído e inquieto. Su
\VUgras{ligereza} y su \VUgras{desobediencia} causan \VUgras{desorden} durante
la clase, y a menudo merece una severa \VUgras{advertencia}. Su
\VUgras{cuaderno de borrador} \textsuperscript{(35)} está sucio, lo llena de
\VUgras{manchas de tinta} con su \VUgras{pluma} \textsuperscript{(34)}, y por eso
siempre usa su \VUgras{goma de borrar} \textsuperscript{(36)}; su \VUgras{cartera}
\textsuperscript{(33)} está rota, porque es muy descuidado. ¿Por qué trae su
\VUgras{portalápices} \textsuperscript{(37)} a la clase de lenguas modernas?

%%K t01-11-2 p
Su vecino Edmundo \textsuperscript{(39)} no lo quiere. Es cierto que es algo
perezoso, pero es callado y se está quieto. No charla; solo que, en vez
de atender, prefiere leer los \VUgras{cuentos} de su \VUgras{libro de
lectura}.

%%K t01-11-3 p
El tercer alumno de este banco es \VUgras{Ludovico} \textsuperscript{(40)}. Es
aplicado. Escribe en una \VUgras{hoja de papel} blanca. Delante de él ha
dejado su \VUgras{papel secante} \textsuperscript{(41)}.

%%K t01-12-1 p
12. --- Los alumnos del segundo banco, a la izquierda del profesor, son
muy jóvenes. El primero, el pequeño \VUgras{Emilio}, todavía no sabe
escribir bien; trae su \VUgras{pizarrín} \textsuperscript{(42)}, su
\VUgras{lápiz de pizarra} y su \VUgras{esponja} \textsuperscript{(43)}.

%%K t01-12-2 p
A su lado están \VUgras{Augusto} y \VUgras{Bernardo} \textsuperscript{(44)}. Este
último trabaja bien, pero su \VUgras{aspecto} es muy descuidado.

%%K t01-13-1 p
13. --- Detrás de ellos hay tres \VUgras{holgazanes}. \VUgras{Alberto} no se
aprende las lecciones. \VUgras{Jacobo} \textsuperscript{(47)} duerme en vez de
escuchar. Su \VUgras{pereza} le acarrea \VUgras{reprimendas} e incluso
\VUgras{castigos}. Por último, \VUgras{Cristián} \textsuperscript{(48)} es demasiado
frívolo. Se \VUgras{porta} mal y no hace ningún esfuerzo por enmendarse.

%%K t01-14-1 p
14. --- Sus vecinos, en cambio, se portan bien. \VUgras{Ernesto}
\textsuperscript{(51)} ama el orden y la regularidad; siempre usa su \VUgras{regla}
\textsuperscript{(50)} y su \VUgras{lápiz} \textsuperscript{(49)}. Es \VUgras{externo}.

%%K t01-14-2 p
\VUgras{Francisco} \textsuperscript{(52)} es \VUgras{interno}; lleva el \VUgras{uniforme}
del instituto; allí come y allí duerme. Es un buen \VUgras{compañero}.

%%K t01-14-3 p
Mauricio afila su \VUgras{lápiz} con su \VUgras{navaja} \textsuperscript{(56)}; es muy
cuidadoso.

%%K t01-14-4 p
Trae todos sus libros, su \VUgras{gramática} \textsuperscript{(55)}, su
\VUgras{cuaderno de notas} \textsuperscript{(54)} y hasta una \VUgras{carpeta de
escritorio} \textsuperscript{(53)}. Su \VUgras{cartera de colegio} \textsuperscript{(57)}
está en el suelo detrás de él.

%%K t01-14-5 p
Las dos mesas del fondo del aula las ocupan \VUgras{buenos alumnos}
\textsuperscript{(19)}.

%%K t01-tit-6 sub
\VUsaut{4.6mm}
\VUpk{10.2pt}{40}{Dibujo y música.}
\VUsaut{3.4mm}

%%K t01-15-1 p
15. --- En el \VUgras{aula de dibujo} hay bonitos \VUgras{modelos} colgados de
la pared y una \VUgras{lámpara} \textsuperscript{(62)} con \VUgras{pantalla},
colgada del techo. El \VUgras{profesor} \textsuperscript{(60)} es muy paciente;
corrige los \VUgras{dibujos} \textsuperscript{(61)} de los alumnos y les enseña a
dibujar un \VUgras{busto} \textsuperscript{(59)} del natural.

%%K t01-16-1 p
16. --- El \VUgras{aula de música} parece muy interesante. Allí se aprende a
leer la \VUgras{música}, a solfear las \VUgras{notas de la escala}
\textsuperscript{(67)} escritas en el \VUgras{pentagrama} (do, re, mi, fa, sol, la,
si\,=\,C. D. E. F. G. A. B.), a reconocer las \VUgras{claves} y a cantar
bonitas \VUgras{melodías}. Las partituras se ponen en el \VUgras{atril}
\textsuperscript{(65)}. Uno de los alumnos \VUgras{toca el piano} \textsuperscript{(64)} y
el \VUgras{profesor de música} \textsuperscript{(66)} \VUgras{lleva el compás}
\textsuperscript{(68)}.

%%K t01-17-1 p
17. --- Alcanzamos a ver un rincón del \VUgras{taller} de \VUgras{trabajos
manuales} \textsuperscript{(69)}, donde los niños pueden aprender a cepillar la
madera y a limar el hierro. No lo hay en todos los institutos.

%%K t01-tit-7 sub
\VUsaut{5.0mm}
\VUpk{10.2pt}{40}{El aula de ciencias.}
\VUsaut{3.6mm}

%%K t01-18-1 p
18. --- El profesor de matemáticas enseña a los alumnos \VUgras{geometría,}
\VUgras{aritmética, cálculo} y el \VUgras{sistema métrico}. Vemos en la
pizarra las principales figuras de geometría: un \VUgras{círculo}
\textsuperscript{(a)}, un \VUgras{cuadrado} \textsuperscript{(b)}, un \VUgras{triángulo}
\textsuperscript{(c)}, una \VUgras{línea} \textsuperscript{(d)}.

%%K t01-18-2 p
Vemos también una \VUgras{operación}; no es una \VUgras{suma}, ni una
\VUgras{resta}, ni una \VUgras{multiplicación}: es una \VUgras{división}
\textsuperscript{(e)}.

%%K t01-19-1 p
19. --- En la lámina del sistema métrico distinguimos: un \VUgras{metro}
\textsuperscript{(a)}, una \VUgras{balanza} \textsuperscript{(b)} y sus \VUgras{pesas}, un
\VUgras{litro} \textsuperscript{(e)}, un \VUgras{decalitro} \textsuperscript{(f)}, un
\VUgras{decilitro} y un \VUgras{metro cúbico} \textsuperscript{(d)}.

%%K t01-19-2 p
Los alumnos estudian también \VUgras{ciencias naturales} \textsuperscript{(11)} ---
zoología \textsuperscript{(a)}, botánica \textsuperscript{(b)}, geología \textsuperscript{(c)} ---
e \VUgras{historia} \textsuperscript{(12)}.

%%K t01-19-3 p
En el armario están los instrumentos de \VUgras{física} \textsuperscript{(15)} y de
\VUgras{química} \textsuperscript{(16)}.

%%K t01-19-4 p
Encima del armario hay: una \VUgras{esfera} \textsuperscript{(14)}, un \VUgras{cono} y
un \VUgras{globo terráqueo} \textsuperscript{(13)}.

%%K t01-19-5 p
Encima de las láminas cuelga un \VUgras{mapa} que representa las cinco
partes del mundo: \VUgras{Europa} \textsuperscript{(g)}, \VUgras{Asia} \textsuperscript{(e)},
\VUgras{África} \textsuperscript{(f)}, \VUgras{América} \textsuperscript{(ab)} y
\VUgras{Oceanía} \textsuperscript{(h)}, y los dos grandes océanos, el
\VUgras{Pacífico} \textsuperscript{(c)} y el \VUgras{Atlántico} \textsuperscript{(d)}.
\VUsaut{6.0mm}
\VUfilet{22mm}
